% Proposed replacement block for the original UCB-AA analysis.
% Root agent will inline this into the single-file main manuscript.
% Complete proofs: theory/ucbaa_repaired_analysis.tex.
\subsection{UCB-AA under a physical sampling budget}
\label{sec:ucbaa_fixed_budget}

We retain the screening architecture of UCB for Arriving Arms (UCB-AA).
At the beginning of round $m$, admit the buffered batch $A_m$ and form
$B_m=C_{m-1}\cup A_m$, where $C_{m-1}$ contains the surviving incumbents.
For $m\ge2$, screen the new arms in phases $p=0,\ldots,m-2$ against
the incumbents and the remaining new arms; only new arms can be removed.
Then apply formal elimination to the merged set in phases $m-1,m$.
Round zero uses phase zero, and round one uses formal phases zero and one.
Empty preliminary batches are no-ops. This retains the distinction between
pre-elimination, which addresses arrival information discrepancy, and
formal elimination after new arms have entered the main competition.

To give this procedure a fixed-budget guarantee, specify its confidence
calibration, comparison data, and stopping rule. For a known physical
budget $T\ge1$ and $0<\delta<1$, let
\begin{equation}
 M_T=\left\lceil\tfrac12\log_2(T+1)\right\rceil,\qquad
 N=8T4^{M_T}/\delta,\qquad\widetilde\Delta_p=2^{-p}.
 \label{eq:ucbaa_main_calibration}
\end{equation}
Retain the original target and radius formulas
\begin{equation}
 n_p=\left\lceil
       \frac{2\log(N\widetilde\Delta_p^2)}{\widetilde\Delta_p^2}
      \right\rceil,\qquad
 b_p=\sqrt{\frac{\log(N\widetilde\Delta_p^2)}{2n_p}}.
 \label{eq:ucbaa_main_targets}
\end{equation}
In an elimination call, sample each candidate to $n_p$ lifetime pulls,
retaining all earlier observations. After the call's sampling is complete,
compare the actual prefix intervals
$\hat r_i(n_p)\pm b_p$. Remove a candidate only if its upper bound is
strictly below the largest reference lower bound. Incumbents have already
reached the necessary prefix before a preliminary comparison, so no
unobserved incumbent estimate is used. Stop after physical pull $T$,
including when it cuts through a target or a round. The calibrated
$n_{M_T}>T$ ensures that the finite phase range covers this execution.
These conventions repair the confidence and budget bookkeeping while
preserving the two original elimination stages. The guarantee below uses
this calibration rather than the smaller empirical calibration in the
submitted manuscript.

\subsection{Roundwise regret, historical loss, and arrival delay}
\label{sec:ucbaa_three_contributions}

Let $s_m$ be the physical start of round $m$,
$A(m)=\bigcup_{j\le m}A_j$, and
$r_{*(m)}=\max_{i\in B_m}r_i$, as in the original roundwise formulation.
Write $\Delta_i^{\mathrm{rd}}(m)=r_{*(m)}-r_i$ for its roundwise gap,
reserving $\Delta_i(t)=r^*(t)-r_i$ for physical time.
Write $r_{\mathrm{adm}}(m)=\max_{i\in A(m)}r_i$ for the best mean
among all admitted arms, including those previously eliminated.
At a physical time $t$, the current round is $m(t)$. The exact regret
identity is
\begin{align}
 R_T(\pi)&=R_{\mathrm{round}}(T)+\mathcal H_T+D_T,\label{eq:ucbaa_main_decomposition}\\
 R_{\mathrm{round}}(T)
  &=\sum_{t\le T}\{r_{*(m(t))}-r_{I_t}\},\nonumber\\
 \mathcal H_T&=\sum_{t\le T}\{r_{\mathrm{adm}}(m(t))-r_{*(m(t))}\},\nonumber\\
 D_T&=\sum_{t\le T}\{r^*(t)-r_{\mathrm{adm}}(m(t))\}.\nonumber
\end{align}
All three quantities are nonnegative. The historical term $\mathcal H_T$ records
the consequence of eliminating an admitted benchmark arm; $D_T$ records
only the time before arrivals are admitted. Combining them into a residual
does not make that residual arrival delay.

For a complete run of $M+1$ rounds, the original $T_0(\pi)$ denotes its
actual random number of pulls, and
$R_{\mathrm{round}}(T_0(\pi))=R_M(\pi)$.
The pathwise identity also applies at this endpoint, but the fixed-budget
bound does not replace $T$ by $\mathbb E T_0(\pi)$ or assert a guarantee
for an arbitrary stopping time. The original $K(m)$, if supplied, bounds
$|A(m)|$; it differs from the actual count $K_t=|\mathcal A_t|$ and is
not an additional input required by the present calibration.

The delay contribution is explicit. Let
$\beta_u=r^*(u)-r^*(u-1)$ for $u\ge2$ and let $\alpha_T(u)$ be the
first round start at or after $u$, with value $T+1$ when no such start
occurs by $T$. Because $A(m)=\mathcal A_{s_m}$,
\begin{equation}
 D_T=\sum_{u=2}^T\beta_u\{\alpha_T(u)-u\}.
 \label{eq:ucbaa_main_delay}
\end{equation}
This formula retains benchmark improvements in the last unfinished round.

\begin{theorem}[A repaired fixed-budget bound for UCB-AA]
\label{thm:ucbaa_main_fixed_budget}
For independent arm streams, each consisting of iid $[0,1]$-valued rewards,
deterministic permanent arrivals, a nonempty initial arm set, and
$K_T=|\mathcal A_T|<\infty$,
the UCB-AA procedure specified
above satisfies, with probability at least $1-\delta$,
\[
 \mathcal H_T=0,\qquad
 R_T(\pi)\le\min\{T,\ 8\sqrt{(2\log N+1)K_TT}+D_T\}.
\]
It follows that
\begin{equation}
 \mathbb E R_T(\pi)\le
 \min\{T,\ 8\sqrt{(2\log N+1)K_TT}
                  +\mathbb E D_T+\delta T\}.
 \label{eq:ucbaa_main_bound}
\end{equation}
The inequalities include incomplete targets and rounds and allow ties
between arm means. No terminal-gap stability condition is needed.
\end{theorem}

The proof uses the actual lifetime observation count. On one simultaneous
prefix-coverage event, a best admitted arm survives every comparison.
A new arm that receives phase-$p$ samples for $p\ge1$ has survived the
preceding comparison in the same round; its roundwise gap is at most
$4b_{p-1}$. An incumbent receives no new samples in formal phase $m-1$,
because it already has $n_{m-1}$ observations. Its new phase-$m$ samples
follow the current round's phase-$m-1$ comparison and obey the same bound.
Thus the $n$th observation of an arm has roundwise loss at most
$4\sqrt{(2\log N+1)/n}$. Summing these distinct prefix charges yields
\eqref{eq:ucbaa_main_bound}; no $G_3$ terminal-gap compression or
union over random elimination chains is used. The supplementary material
gives the full proofs.

For $T\ge2$, the choice $\delta=T^{-2}$ makes the statistical term
$O(\sqrt{K_TT\log(T+1)})$. Sublinear regret therefore follows under
the sufficient conditions $K_T\log(T+1)=o(T)$ and
$\mathbb E D_T=o(T)$. The delay condition is essential to this statement;
the result does not assert sublinear regret for arbitrary arrivals.
For an unknown horizon, a restart at physical budgets $1,2,4,\ldots$
gives a valid extension with the same statistical order and its own explicit
delay term. The supplementary material provides its construction.
This physical-time wrapper is distinct from the withdrawn argument based
on completing geometrically many algorithmic rounds.

\begin{corollary}[Static arms and fixed-gap distractors]
\label{cor:ucbaa_main_fixed_gap}
Suppose an arm available at time one has mean $r^\star$ and every arm
arriving by $T$ has mean at most $r^\star$. Write
$\Delta_i=r^\star-r_i$ and $A=2\log N+1$. Then $D_T=0$ pathwise and,
with probability at least $1-\delta$,
\begin{equation}
 R_T(\pi)\le\sum_{\substack{i\in\mathcal A_T\\\Delta_i>0}}
       \min\left\{\Delta_i(T-\tau_i+1),\frac{16A}{\Delta_i}\right\}.
 \label{eq:ucbaa_main_fixed_gap}
\end{equation}
Adding $\delta T$ gives an expected-regret bound, and both bounds may
be capped at $T$.
\end{corollary}

Indeed the active best mean equals $r^\star$ on the common confidence
event. Applying the actual-prefix charge to a suboptimal arm's last
observation gives $N_i(T)\le16A/\Delta_i^2$, while
$N_i(T)\le T-\tau_i+1$. For $T\ge2$ and $\delta=T^{-2}$, the static case therefore
has upper order $\sum_{i:\Delta_i>0}\log(T+1)/\Delta_i$; arrivals with
a fixed positive gap have upper order $K_T\log(T+1)$.
For example, retain the original illustration with an initial optimum
$r^\star=3/4$, another initial mean $1/4$, and arriving means
$r_i=1/2-2^{-i-2}$ at exogenous physical times $\tau_i$, $i\ge1$.
All distractors have gap at least $1/4$. This preserves the valid
fixed-gap conclusion without the withdrawn $G_3$ simplification or a
claim that removing preliminary screening forces linear regret.
